\documentclass[11pt]{article}

% Change "review" to "final" to generate the final (sometimes called camera-ready) version.
% Change to "preprint" to generate a non-anonymous version with page numbers.
\usepackage[]{acl}
\usepackage{times}
\usepackage{latexsym}
\usepackage{amsmath}
\usepackage{amssymb}
\usepackage{xurl}

% For proper rendering and hyphenation of words containing Latin characters (including in bib files)
\usepackage[T1]{fontenc}
% This assumes your files are encoded as UTF8
\usepackage[utf8]{inputenc}
\usepackage{microtype}
\usepackage{inconsolata}
\usepackage{graphicx}

\title{Group 90 Progress Report:\\Sleep Debt and Screen Time: Late Night Phone Habits}

\author{
  Nadeem Elsayed \and Muhammad Ullah \\
  Department of Computing and Software \\
  McMaster University \\
  \texttt{\{elsayn3, ullahm7\}@mcmaster.ca}
}

\begin{document}
\maketitle

\section{Introduction}
\label{sec:intro}

Smartphones are an integral part of our lives due to benefits; such as easy accessibility to information, social connectivity, workplace/financial applications, and entertainment including watching videos or playing video games. Smartphones have been revolutionized the ease of information and has everyone connected. However, the smartphone - with many pros stated - has also many drawbacks that has disconnected everyone from reality. Spending too much time on the screen reduced personal socialization and significant sleep \citep{cbc2026}. Sleep plays a crucial role in developing new memory and is essential for all stages of learning. With inadequate amounts of sleep, the body and mind experience rapid physical and cognitive decline, affecting overall long-term health. These include mood swings, \citep{shah2024effects}. Over 3/4 of Americans lose sleep because of digital distractions \citep{newswise2026}. The number might be different for different nations, but this is a severe problem that affects a lot of people. 

The goal of this project is to analyze and predict the fatigue that people feel the following day based on their screentime usage. 

There are a few challenges with trying to predict `sleep fatigue' based on screentime. There are a lot of variables involved, they might not scale linearly with sleep loss, and there is a limit to the amount of sleep fatigue a person might experience (it could be different for different people). 

The data type is tabular containing number and category features like screen time, type of screen use, device type, and the resulting sleep patterns during weekdays and weekends. The task type is regression since we are trying to figure out the exact amount of sleep loss in hours. The number of classes is not applicable for this task since we are using regression. Since we are predicting sleep loss for both weekdays and weekends, this is a multi-label problem where we will get 2 outputs per set of inputs.

\section{Related Work}

Similar studies have been done to show the effects of screen usage on sleep. A systematic review of school-aged children and adolescents \citep{hale2015} found that screen time is generally linked to worse sleep. A meta-analysis covering 125,198 children (mean age 14.5) reported that bedtime media-device use was associated with an odds ratio of 2.17 for inadequate sleep quantity \citep{carter2016}. These studies report associations but do not predict an individual student’s sleep loss. Their findings support using screen-use variables as inputs in this study, particularly the days per week each device is used at bedtime and the daily hours of each screen activity, and they suggest that higher screen use should be linked to greater sleep loss.

However, Alam et al. \citet{alam2023} used objectively recorded smartphone usage from 269 participants aged 15–64 and found, with correlation and regression analyses, that more use was associated with shorter sleep and later bedtimes. This supports treating sleep duration as a numeric outcome that can be modelled from screen usage. It also highlights a limitation of the present dataset, where screen time is self-reported rather than measured. Simon et al. \citet{simon2024} trained six algorithms, including SVM, XGBoost and random forest, to detect insomnia symptoms from smartphone usage and reported low sensitivity (0.05–0.27) on the test subsample, a warning that predicting sleep from behaviour is hard. This sets realistic expectations of only moderate performance and shows why every model must be compared against simple baselines on held-out data.

The most similar study that was found, a 2024 conference abstract in SLEEP \citep{sleep2024abstract}, used LASSO and random forest to rank factors associated with adolescent sleep outcomes, including sleep duration. It ranks factors rather than predicting hours of sleep loss. Its approach supports the modelling choices here: ridge regularization is applied to polynomial features to keep all inputs while shrinking noisy coefficients, which suits the many overlapping screen-use variables, and lasso regularization is planned as a next step to identify which variables are unnecessary. A public GitHub project \citep{som2024} classifies sleep quality from screen exposure with logistic regression, random forest and SVM but treats it as a classification problem. Its use of tree-based models indicates that non-linear models such as decision trees are a reasonable comparison, and these are also planned.

No study was found that predicts weekday and weekend sleep loss in hours for school-aged adolescents. The present study addresses this gap by combining the predictive framing of Alam et al. \citet{alam2023} and Simon et al. \citet{simon2024} with the screen-use variables identified by Hale and Guan \citet{hale2015} and Carter et al. \citet{carter2016}.
\section{Dataset and Preprocessing}
\label{sec:dataset}

\subsection{Data Source and Characteristics}
We're using the research dataset by \citet{perez2022screen}, hosted on the Mendeley Data repository under a Creative Commons open-access license. The survey records data collected across 52 secondary schools in Argentina. The survey captures metrics such as demographic attributes like age and gender, device access, the types of digital screentime (gaming, social media, and video streaming), sleep timing, and clinical screening tools to look for narcolepsy and other disorders.

The raw archive (\texttt{repositorio.xlsx}) contains 4,115 raw respondent submissions with 123 distinct tabular variables recorded in Spanish.

\subsection{Data Cleaning and Feature Harmonization}
To establish a machine learning pipeline, we first implemented a data translation pipeline:
\begin{enumerate}
    \item \textbf{Translation and Feature Extraction:} All feature headings and the responses were mapped from Spanish to English (ex: tpopantallatotalhs -> screen\_total\_hrs, tporedeshs -> screen\_social\_media\_hrs, andtposueñosem -> sleep\_weekday\_hrs).
    \item \textbf{Cohort Filtering and Missingness:} We isolated the participants from 11 to 20 years. Records with missing values for primary sleep durations or total screen hours were eliminated. Sub-activity values (such as gaming or console hours) where respondents indicated no use were set to 0 hours.
    \item \textbf{Clinical Plausibility Bounding:} Outlier responses that did not make physical sense for humans were excluded (restricting sleep durations to $3.0 \le \text{sleep} \le 16.0$ hours/day and daily screen exposure to $0.0 \le \text{screen} \le 24.0$ hours/day).
\end{enumerate}
After this preprocessing, our finalized dataset is composed of 2,496 student records, which have been exported to \./dataset\_clean\_english.csv.

\subsection{Target Formulation: Sleep Debt}
Consistent with pediatrician recommendations made by \citet{paruthi2016recommended}, the target sleep duration for adolescents is set to $S_{\text{rec}} = 9.0$ hours. We define the target vector $\mathbf{y} = [y_{\text{weekday}}; y_{\text{weekend}}] \in \mathbb{R}^2$ as:
\begin{align}
    y_{\text{weekday}} &= 9.0 - T_{\text{sleep, weekday}} \\
    y_{\text{weekend}} &= 9.0 - T_{\text{sleep, weekend}}
\end{align}
Under this formula, positive values represent severe sleep loss (sleep debt), whereas negative values represent sleep surplus. In our preprocessed data, adolescents experience an average weekday sleep loss of $\mu = 1.96$ hours ($\sigma = 1.36$ hours), contrasted with a weekend sleep surplus of $\mu = -0.42$ hours ($\sigma = 1.91$ hours), with a mean social jetlag of 2.38 hours of rebound sleep on weekends.

\section{Features}

The model inputs are 13 numeric features selected from the 123 columns of the raw dataset: age, sex (coded 0 for female and 1 for male), body mass index (BMI), total daily screen time, daily hours spent on four screen activities (social media, video games, games on a phone or tablet, and TV or online video), and the number of days per week that each of five devices (cell phone, computer, tablet, game console and TV) is used. The two sleep durations are not inputs. They are used only to compute the targets, weekday and weekend sleep loss.

Feature selection was performed manually. Variables unrelated to sleep and screen time (such as school performance and subjective fatigue responses) and variables that were hard to interpret or redundant were removed, so that the model is not crowded with parameters unlikely to affect the predicted outcome. Variables derived from the targets, such as social jetlag (weekend minus weekday sleep), were also excluded because they would leak the answer. Total screen time is correlated with the activity-level hours, since it includes them, and ridge regularization is used to handle this overlap.

Feature engineering consisted of three steps. First, column names were translated and renamed. Second, missing activity-level hours and device-use days were filled with 0, interpreted as no use. Third, a degree-2 polynomial expansion was applied to the 13 inputs, producing 104 features: 13 linear terms, 13 squared terms and 78 pairwise interaction terms. This allows the model to capture non-linear effects and interactions, such as the combined effect of social media hours and phone use. The inputs are standardized (mean 0, standard deviation 1) before the expansion, using statistics from the training set only, so that the ridge penalty treats all terms equally. No embeddings or representation learning are used, and the model is not a neural network; the feature vectors are passed directly to scikit-learn regressors. Lasso regularization is planned as a data-driven form of feature selection, to identify which variables and interaction terms are unnecessary.

\section{Model Implementation}
\label{sec:impl}

We represent our task as a multi-target linear and non-linear regression problem where we learn a parameter mapping $f: \mathbb{R}^{13} \to \mathbb{R}^2$ optimizing the Mean Squared Error (MSE):
\begin{equation}
    MSE(W) = \frac{1}{N} \sum_{i=1}^N \|\mathbf{y}_i - \hat{\mathbf{y}}_i\|_2^2
\end{equation}

To establish benchmarks and evaluate model learning, we implemented three candidate solutions with scikit-learn:

\subsection{Baseline 1: Empirical Mean Regressor}
We construct a naive baseline predicting the training set mean vector for every test instance:
\begin{equation}
    \hat{\mathbf{y}}_{\text{dummy}} = \bar{\mathbf{y}}_{\text{train}} = \frac{1}{N_{\text{train}}} \sum_{i=1}^{N_{\text{train}}} \mathbf{y}_i
\end{equation}
This baseline represents the lower bound ($R^2 \approx 0.0$).

\subsection{Baseline 2: Multiple Linear Regression (OLS)}
We use a multi-output Ordinary Least Squares (OLS) solution:
\begin{equation}
    \hat{\mathbf{y}} = \mathbf{W}^T \tilde{\mathbf{x}} + \mathbf{b}
\end{equation}

\subsection{Candidate Model: Polynomial Regression with Ridge Regularization}
Recognizing that screen exposure exhibits non-linear effects, we construct a degree-2 polynomial feature mapping $\phi(\mathbf{x})$ encompassing all linear, quadratic, and pairwise interaction terms:
\begin{equation}
    \phi(\mathbf{x}) = [1, x_1, \dots, x_d, x_1^2, x_1 x_2, \dots, x_d^2]^T \in \mathbb{R}^{104}
\end{equation}
To prevent overfitting on the 104-dimensional space, we use the Ridge objective with penalty parameter $\alpha = 100.0$:

\begin{equation}
    {L}_{\text{Ridge}} = \text{MSE} + \alpha \|\mathbf{w}\|_2^2
\end{equation}

\begin{table*}[t]
\centering
\small
\begin{tabular}{lcccccc}
\hline
\textbf{Model Architecture} & \multicolumn{3}{c}{\textbf{Weekday Sleep Loss (Hours)}} & \multicolumn{3}{c}{\textbf{Weekend Sleep Loss (Hours)}} \\
\cline{2-4} \cline{5-7}
& \textbf{MAE} $\downarrow$ & \textbf{RMSE} $\downarrow$ & $\mathbf{R^2}$ $\uparrow$ & \textbf{MAE} $\downarrow$ & \textbf{RMSE} $\downarrow$ & $\mathbf{R^2}$ $\uparrow$ \\
\hline
Baseline 1: Mean Regressor & 1.043 & 1.406 & -0.000 & 1.549 & 1.949 & -0.002 \\
Baseline 2: Linear Regression (OLS) & \textbf{0.971} & \textbf{1.332} & \textbf{0.103} & \textbf{1.511} & \textbf{1.920} & \textbf{0.027} \\
Candidate: Polynomial Ridge (Deg 2, $\alpha=100$) & 0.988 & 1.360 & 0.065 & 1.535 & 1.967 & -0.021 \\
\hline
\end{tabular}
\caption{Preliminary quantitative performance across 500 held-out test adolescents. Metrics report Mean Absolute Error (MAE in hours), Root Mean Squared Error (RMSE in hours), and Coefficient of Determination ($R^2$).}
\label{tab:results}
\end{table*}

\section{Results and Evaluation}
\label{sec:results}

\subsection{Evaluation Protocol}
We split the preprocessed dataset N = 2,496 using an 80/20 train-test split with fixed random seed ($N_{\text{train}} = 1{,}996$, $N_{\text{test}} = 500$). Performance is measured across three standard regression metrics:
\begin{align}
    \text{MAE} &= \frac{1}{N_{\text{test}}} \sum_{i=1}^{N_{\text{test}}} |y_i - \hat{y}_i| \\
    \text{RMSE} &= \sqrt{\frac{1}{N_{\text{test}}} \sum_{i=1}^{N_{\text{test}}} (y_i - \hat{y}_i)^2} \\
    R^2 &= 1 - \frac{\sum_{i} (y_i - \hat{y}_i)^2}{\sum_i (y_i - \bar{y})^2}
\end{align}

\subsection{Preliminary Findings}
The test results are detailed in Table~\ref{tab:results}. Multiple Linear Regression achieves a Mean Absolute Error of 0.971 hours on weekday sleep loss (R\^2 = 0.103, RMSE = 1.332 hours), demonstrating that it's better than the baseline ($\text{MAE} = 1.043$ hours).

The degree-2 polynomial model with Ridge regularization ($\alpha = 100$) has a slightly higher test MAE ($0.988$ hours) than standard linear regression. This suggests that expanding to 104 quadratic and interaction features adds noise that outweighs non-linear gains when constrained to a linear regularizer. 

Furthermore, all models showed lower predictive accuracy on weekend sleep loss ($R^2 \approx 0.027$) compared to weekday sleep loss ($R^2 \approx 0.103$). Weekend sleep schedules among teenagers show weirder behaviour and compensatory rebound sleeping, which is weakly correlated with isolated screen time.

\section{Feedback and Plans}

Feedback was received by email from two TAs, Abdelaziz Serour and Shreya Adrita Banik and both agreed on the main issues. Abdelaziz said: Investigate why polynomial regression performs worse with and without regularization, tuning regularization strength, try Lasso and examine which feature coefficients become zero and compare its results with Ridge and linear regression, 0\% of the variation in weekday sleep loss and 3\% in weekend sleep loss so investigate whether the available features are sufficiently related to sleep loss, expansion of 13 input features into 104 polynomial features which may introduce unnecessary interactions, and he sees weekend predictions are worse than weekday predictions so analyze which features are important for each target. Shreya had also pointed out that all the current models have relatively low $R^2$ scores, so explore some other nonlinear model. She has said It can be useful to see which features contribute most to the prediction and discuss why the models might be struggling and considers using  cross-validation in the future to see if results are consistent across different splits. 
 
Based on this feedback and a review of the implementation, the remaining work is as follows. First, the cause of the polynomial model’s weaker performance will be investigated by comparing degree-1 and degree-2 models, each with and without regularization, to determine whether the polynomial features or the penalty are responsible. The expansion to 104 features may add noise, so a smaller set of selected interaction terms will also be tested. Second, alpha will be chosen through k-fold cross-validation on the training data instead of being fixed, and the same procedure will be applied to Lasso. Features are already standardized before regularization, which is required for both methods. Lasso will also show which coefficients are reduced to zero, which gives evidence about the noisy features suspected so far. Third, a possible implementation of decision tree or random forest regressor will be added as a non-linear comparison, since every current model has a low R². Fourth, weekday and weekend sleep loss will be modelled separately, and feature importance will be examined for each, because the weekend predictions are weaker and may depend on different factors. 

\section*{Team Contributions}

Nadeem Elsayed developed the data pipeline and models. This included cleaning the raw dataset and exporting the cleaned file, building the sleep loss targets, and implementing the mean regressor, linear regression and polynomial ridge regression baselines. He also ran the evaluation and wrote the Dataset and Preprocessing, Model Implementation, and Results and Evaluation sections, and he coordinated the TA feedback.

Muhammad Ullah wrote the Related Work, Features, and Feedback and Plans sections, worked on Nadeem’s code and results for the Features section, and compiled the TA feedback into the plans for the remaining work. He also set up the LaTeX build and helped proofread and format the report. Both members reviewed the report before submission.

\bibliography{custom}

\end{document}
